\documentclass[10pt,a4paper]{amsart}
\usepackage{amsmath,amssymb,amsthm,graphicx}
\usepackage{amsmath}
\usepackage{amssymb}
\usepackage{hyperref}
\newtheorem{Theorem}{Theorem}[section]
\newtheorem{Corollary}[Theorem]{Corollary}
\newtheorem{Lemma}[Theorem]{Lemma}
\newtheorem{Proposition}[Theorem]{Proposition}
{\theoremstyle{definition}
\newtheorem{Note}[Theorem]{Note}
\newtheorem{Definition}[Theorem]{Definition}
\newtheorem{Problem}[Theorem]{Problem}
}
\usepackage{cleveref}
\crefname{originaldefinition}{original definition}{original definitions}
\Crefname{originaldefinition}{Original definition}{Original definitions}
\makeatletter
\newlabel{def:abc}{{1.3}{source}{}{sourcepointer-1}{}}
\newlabel{def:abc@cref}{{[originaldefinition][0][]1.3}{[1][0][]source}{}{}{}}
\makeatother
\crefname{originaldefinition}{original definition}{original definitions}
\Crefname{originaldefinition}{Original definition}{Original definitions}
\makeatletter
\newlabel{def:uaw}{{1.2}{source}{}{sourcepointer-2}{}}
\newlabel{def:uaw@cref}{{[originaldefinition][0][]1.2}{[1][0][]source}{}{}{}}
\makeatother
\crefname{originallemma}{original lemma}{original lemmas}
\Crefname{originallemma}{Original lemma}{Original lemmas}
\makeatletter
\newlabel{lem:multtable2}{{1.12}{source}{}{sourcepointer-3}{}}
\newlabel{lem:multtable2@cref}{{[originallemma][0][]1.12}{[1][0][]source}{}{}{}}
\makeatother
\crefname{originallemma}{original lemma}{original lemmas}
\Crefname{originallemma}{Original lemma}{Original lemmas}
\makeatletter
\newlabel{lem:xyz}{{1.11}{source}{}{sourcepointer-4}{}}
\newlabel{lem:xyz@cref}{{[originallemma][0][]1.11}{[1][0][]source}{}{}{}}
\makeatother
\title{The universal Askey-Wilson algebra maps onto the Askey-Wilson algebra, sending alpha, beta, gamma to a fixed central scalar}
\author{Paul Terwilliger}
\date{doi:10.3842/SIGMA.2011.069; CC BY4.0}
\begin{document}
\maketitle
\noindent\textit{Editorial note.} All delivered mathematical text is the author's own, copied from the pinned publisher source sources/sigma-2011-069-askey-wilson/pinned-original-source.tex; no statement or proof step was rewritten, reordered or removed.\par\medskip
\noindent\textit{Editorial note.} Preamble provenance: the retained preamble is the source's own layout adapter, if any, plus the author's macro block at lines 10-18 of sources/sigma-2011-069-askey-wilson/pinned-original-source.tex (everything between the publisher class metadata and \textbackslash{}begin\{document\}, comments excluded). Nothing was added, no macro was invented, and the publisher document class is deliberately not reproduced; the excerpt serves only as preamble, never as a span coordinate.\par\medskip
\noindent\textit{Editorial note.} External references are navigation pointers only; every deliberately kept step of the author's argument stays inside the counted spans. The displayed numbers of those pointers are the author's own numbering under the author's \textbackslash{}newtheorem declarations (the source has no subsections, so numbering is section.index), computed mechanically by the lane's scanner.\par\medskip
\noindent\textit{Editorial note.} Four references to other labelled items occur inside the counted spans (Lemmas lem:xyz and lem:multtable2, Definitions def:uaw and def:abc); each is recorded once as a navigation stub with the matching kind and the author's own number, and every step that consumes them stays inside the counted span.\par\medskip
\noindent\textit{Editorial note.} The counted spans contain no bibliography citation, so no bibliography excerpt is attached to this card.\par\medskip
\noindent\textit{Allowed standard prerequisites (editorial).}\begin{itemize}
\item The evaluation identities for the generators A, B, C of the Askey-Wilson algebra Lambda. Same source, Lemmas labelled lem:xyz and lem:multtable2; invoked inside the counted proof at lines 948-964. Supply the three displayed identities satisfied by the images A\textasciicircum{}vee, B\textasciicircum{}vee, C\textasciicircum{}vee.
\item The defining relations of the universal Askey-Wilson algebra. Same source, Definition labelled def:uaw; invoked inside the counted proof at lines 965-971. Turns the three identities into the relations that guarantee existence of the homomorphism.
\item The definitions of alpha, beta, gamma and of the scalar mu. Same source, Definition labelled def:abc; invoked inside the counted proof at lines 974-978. Identifies the images of alpha, beta, gamma with nu I, where nu = (q\textasciicircum{}2+q\textasciicircum{}\{-2\}) mu + mu\textasciicircum{}2.
\end{itemize}\medskip
\ifdefined\hypertarget\hypertarget{sourcepointer-1}{}\fi
\noindent\textit{Source reference.} Original definition 1.3: sources/sigma-2011-069-askey-wilson/pinned-original-source.tex, source lines 977--977. Navigation pointer for \textbackslash{}ref\{def:abc\} used inside the counted span (line 977); the argument that consumes it is retained in the span.\par\medskip
\ifdefined\hypertarget\hypertarget{sourcepointer-2}{}\fi
\noindent\textit{Source reference.} Original definition 1.2: sources/sigma-2011-069-askey-wilson/pinned-original-source.tex, source lines 971--971. Navigation pointer for \textbackslash{}ref\{def:uaw\} used inside the counted span (line 971); the argument that consumes it is retained in the span.\par\medskip
\ifdefined\hypertarget\hypertarget{sourcepointer-3}{}\fi
\noindent\textit{Source reference.} Original lemma 1.12: sources/sigma-2011-069-askey-wilson/pinned-original-source.tex, source lines 949--949. Navigation pointer for \textbackslash{}ref\{lem:multtable2\} used inside the counted span (line 949); the argument that consumes it is retained in the span.\par\medskip
\ifdefined\hypertarget\hypertarget{sourcepointer-4}{}\fi
\noindent\textit{Source reference.} Original lemma 1.11: sources/sigma-2011-069-askey-wilson/pinned-original-source.tex, source lines 948--948. Navigation pointer for \textbackslash{}ref\{lem:xyz\} used inside the counted span (line 948); the argument that consumes it is retained in the span.\par\medskip
\begin{Lemma}
\label{prop:uhomoh}
There exists a unique $\mathbb F$-algebra homomorphism
$\pi:\Delta \to \Lambda $
 that sends
\begin{gather*}
A \mapsto q \mathcal A+q^{-1}{\mathcal A}^{-1},
 \qquad
B \mapsto q \mathcal B+q^{-1}{\mathcal B}^{-1},
\qquad
C \mapsto q\mathcal C+q^{-1}{\mathcal C}^{-1}.
\end{gather*}
The homomorphism $\pi$ sends
\begin{gather}
 \alpha \mapsto \nu I,
 \qquad
\beta \mapsto  \nu I,
 \qquad
\gamma \mapsto \nu I,
\label{eq:pisend}
\end{gather}
where $\nu= (q^2+q^{-2})\mu + \mu^2$.
\end{Lemma}

\begin{proof}
Def\/ine
\begin{gather*}
A^\vee = q\mathcal A+q^{-1}{\mathcal A}^{-1},
\qquad
B^\vee =q\mathcal B+q^{-1}{\mathcal B}^{-1},
\qquad
C^\vee = q\mathcal C+q^{-1}{\mathcal C}^{-1}.
\end{gather*}
By Lemma~\ref{lem:xyz}
and Lemma~\ref{lem:multtable2},
\begin{gather}
(q+q^{-1})A^\vee+ \frac{qB^\vee C^\vee-q^{-1}C^\vee B^\vee}{q-q^{-1}}
  = \nu I,
\label{eq:vee1}
\\
(q+q^{-1})B^\vee+
\frac{qC^\vee A^\vee -q^{-1}A^\vee C^\vee }{q-q^{-1}}
 =  \nu I,
\label{eq:vee2}
\\
(q+q^{-1})C^\vee +
\frac{qA^\vee B^\vee -q^{-1}B^\vee A^\vee }{q-q^{-1}}
 =  \nu I.
\label{eq:vee3}
\end{gather}
By
(\ref{eq:vee1})--(\ref{eq:vee3}) and
since
 $\nu I$ is central, the elements
$A^\vee$, $B^\vee$, $C^\vee$ satisfy the
def\/ining relations for~$\Delta $ from
Def\/inition~\ref{def:uaw}.
Therefore
the homomorphism $\pi$ exists. The homomorphism $\pi$
is unique since $A$, $B$, $C$ generate~$\Delta $. Line
(\ref{eq:pisend})
follows from
Def\/inition~\ref{def:abc} and
(\ref{eq:vee1})--(\ref{eq:vee3}).
\end{proof}

\begin{thebibliography}{99}

\footnotesize\itemsep=0pt
\bibitem{alpern}
Alperin R.C.,
Notes: ${\rm  {PSL}}_2(\mathbb Z)= \mathbb Z_2 \star \mathbb Z_3$,
\href{http://dx.doi.org/10.2307/2324963}{{\it Amer. Math. Monthly}} {\bf 100} (1993), 385--386.
\end{thebibliography}
\end{document}
